\documentclass[11pt]{article}
\usepackage{icmlko}
\begin{document}

\title{축을 다 빼도 안 죽는데 깊이 하나가 죽인다 --- 하루 만에 고친 읽기}
\author{월드모델 실험실}
\date{노트 521 · 2026년 8월 4일}
\maketitle

\begin{abstract}
노트 520 이 되살아남의 비대칭을 ``\textbf{용량은 무엇을 표현할 수 있나를
정하고 자료 쓰는 법은 어느 행을 보나만 바꾼다}''로 읽고, 따라오는 예측을
사전 등록했다 --- \emph{축을 늘리는 변경도 표현력을 늘리므로 자료 쪽 갈래를
되살릴 수 있어야 한다.} 대우로 시험했다: \textbf{축을 빼면 \texttt{TRAINR}
이득이 줄어야 한다.}

\begin{center}
\begin{tabular}{lrl}
\hline
설정 & \texttt{TRAINR} 판 이득 & 부호 \\
\hline
챔피언(축 다) & $+0.0028$ & $10/12$ \\
\texttt{SPEC} 끔 & $+0.0032$ & $12/12$ \\
\texttt{SEASON} 끔 & $+0.0019$ & $11/12$ \\
\textbf{둘 다 끔} & $\mathbf{+0.0036}$ & $12/12$ \\
\hline
\textbf{깊이 $6$} & $\mathbf{-0.0000}$ & $5/12$ \\
\hline
\end{tabular}
\end{center}

\textbf{반증.} 축을 빼도 이득이 안 줄고 \emph{둘 다} 빼면 오히려 는다
($+0.0036$). \textbf{깊이만 죽인다.}
\end{abstract}

\begin{figure}[t]
\centering
\includegraphics[width=\textwidth]{figs/f1.pdf}
\caption{\textbf{가} --- 다섯 설정. \textbf{나} --- 고친 읽기.
\textbf{다} --- 하루.}
\end{figure}

\section{읽기를 고친다}

\begin{quote}
\textbf{옛 읽기}(노트 520) --- 용량은 \emph{표현력}을 정하고, 표현력이
늘면 자료 쪽 후보가 새로 값을 갖는다.

\textbf{새 읽기}(이 노트) --- 깊이는 \emph{상호작용 차수}를 정한다. 한
갈래에서 몇 개를 조합할 수 있느냐다. \textbf{최근 가중이 사는 것은 깊은
상호작용}이라 깊이 $6$ 에서는 표현이 안 되고, \textbf{열이 몇 개든
상관없다.}
\end{quote}

축은 ``무엇이 있나''만 정한다. 있는 것을 조합하는 능력이 없으면 열을
늘려도 소용이 없고, 조합할 수 있으면 열이 적어도 된다 --- 실제로 축 둘을
빼자 이득이 $+0.0028 \to +0.0036$ 으로 \emph{늘었다}(다른 축이 없으니
남은 축의 깊은 조합에 더 기댄다는 읽기와 맞는다).

\textbf{정직하게} --- \texttt{SEASON} 을 빼면 $+0.0019$ 로 조금 준다.
계절은 \emph{시간} 축이라 최근 가중과 같은 재료를 쓴다. 그것은
``표현력''이 아니라 \textbf{같은 신호를 두 곳에서 쓰는 것}이다.

\section{조항을 다시 좁힌다}

\begin{quote}
\textbf{조항(정정 $2$)} --- \textbf{깊이를 바꾸면 그 깊이에서 닫았던 자료
쪽 갈래를 다시 연다.} ``용량''이나 ``표현력''이 아니라 \textbf{깊이}다.
\texttt{it} $\cdot$ $lr$ $\cdot$ $K$ $\cdot$ 축 개수는 이 성질이 없다
(노트 520 이 \texttt{it} $\cdot$ $lr$ $\cdot$ $K$ 를, 이 노트가 축을
확인했다).
\end{quote}

이틀 사이에 같은 조항을 세 번 고쳤다 --- 노트 519 ``용량이 바뀌면'' →
노트 520 ``자료 쪽만, 역은 아님'' → 이 노트 ``깊이만''. \textbf{매번 좁아진
것이 좋은 신호다}(반대 방향이면 이야기를 늘리고 있다는 뜻이다).

\section*{남은 것}

\begin{itemize}
\item \textbf{문턱이 어디인가} --- 깊이 $6$ 에서 $0$, $12$ 에서 $+0.0028$.
$8 \cdot 10 \cdot 14$ 를 재고 있다. 문턱이 뚜렷하면 ``몇 차 상호작용이
필요한가''를 셀 수 있다.
\item \textbf{다른 자료 후보도 깊이 감응인가} --- \texttt{TRAINW} $\cdot$
\texttt{TRAINCAP} $\cdot$ \texttt{DOMDROP} 을 깊이 $6$ 에서 재면 같은
그림이 나오나. 나오면 \textbf{이 사이클 전체가 노트 457 의 자식}이다.
\item \textbf{깊이 $6$ 시절에 죽은 다른 자료 후보} --- 원장에 더 있다.
\end{itemize}

\end{document}
